\begin{problem}{H}{Helix Havoc}{1 second}

\noindent
Kogoro and Conan were invited to stay overnight at a secluded castle belonging to the wealthy Mawari family. The family was holding a grand reunion that night to celebrate the return of their most precious heirloom, the legendary \textbf{Helix Plate}, a giant rectangular plate forged entirely from pure gold.

\noindent
The plate was said to hold a peculiar secret. Its surface was engraved with the numbers $1,2,3,\ldots$ in an increasing \textbf{counterclockwise} spiral starting from the \textbf{upper-left} corner. The exact dimensions of the plate had never been publicly recorded, but the family knew that \textbf{both dimensions were odd and at least $3$}. The \textbf{Helix Number} of the plate is defined by taking the median of each row and then finding the median of those row medians.

\noindent
For example, when the dimensions are $(N,M)=(3,5)$, the plate is filled as follows:

\begin{tabular}{ccccc}
  1 & 12 & 11 & 10 & 9 \\
  2 & 13 & 14 &15 & 8 \\
  3 & 4 & 5 & 6 & 7
\end{tabular}

\noindent
The medians of the three rows are $10$, $13$, and $5$, respectively. Taking the median of these three values gives $10$. Therefore, the \textbf{Helix Number} of this plate is $10$.

\noindent
Late that night, a loud bang echoed through the castle. Everyone rushed toward the head of the Mawari family's bedroom, only to find the door locked from the inside. After forcing the door open, they discovered the head of the family lying dead on the floor. The safe had been opened, and the Helix Plate was nowhere to be found.

\noindent
Conan immediately noticed that the victim's hand was pointing toward the empty safe, and there was a number $K$ written in blood on the floor next to the body. He suspected that this was the victim's final clue, deducing that $K$ represented the \textbf{Helix Number} of the stolen plate.

\noindent
Determine all possible dimensions $(N,M)$ of the missing Helix Plate whose Helix Number is exactly $K$. Print at most $50$ such pairs.

\InputFile

\noindent
The first line of the input contains a single integer $T$ ($1 \le T \le 1000$) --- the number of testcases.

\noindent
Each of the following $T$ lines contains a single integer $K$ ($1 \le K \le 10^{15}$), the \textbf{Helix Number}.

\noindent
It is guaranteed that the sum of $K^{1/3}$ over all testcases does not exceed $10^5$.

\OutputFile

\noindent
For each testcase, print the answer in the following format:

\noindent
On the first line, print a single integer $C$ --- the total number of valid pairs $(N, M)$ that result in a \textbf{Helix Number} equal to $K$.

\noindent
Then, print $\min(C, 50)$ lines. Each of these lines should contain two space-separated integers $N$ and $M$, representing the dimensions of a valid plate. If there are multiple ways to choose $50$ pairs, any valid subset will be accepted.

\Examples

\begin{example}
\exmp{1}{3
10
14
22
}{1
3 5
0
2
3 13
7 5
}
\end{example}

\pagebreak

\end{problem}